\section{Related Work}\label{sec:related}

\textbf{Parallel model checking of untimed systems.}
Parallel explicit-state BFS is well studied~\cite{Stern01}.  The standard
approach uses a globally shared hash table with fine-grained locking or
lock-free operations (compare-and-swap).  Our setting differs in two ways:
(i)~symbolic states carry a zone, not just a configuration, making hashing and
equality non-trivial; (ii)~subsumption (inclusion between zones) means that
the semantics of ``already seen'' is richer than equality.  These two
properties make concurrent passed-list updates significantly harder to reason
about, which is why we opt for a sequential merge instead.

\textbf{Parallel symbolic model checking of timed systems.}
\textsc{LTSmin}~\cite{Kant15} parallelises timed-automata exploration using a
global hash table, but it stores one entry per (configuration, zone) pair
\emph{without} subsumption.  As a result it stores more symbolic states than
a subsumption-based tool and explores more transitions.  Our approach retains
subsumption and therefore explores the same (minimal) number of states as
sequential UPPAAL.

\textbf{Level-synchronous BFS for model checking.}
The ``frontier'' structure (sometimes called BFS-level parallelism) appears
in distributed-memory model checkers such as \textsc{DiVinE}~\cite{Barnat12}
and in GPU-accelerated BFS~\cite{Merrill12}.  The technical novelty here is
applying it to symbolic zone-graph exploration, where the independence of
$\succfn$ calls within a level is less obvious because states share the same
data structures (passed list), and arguing correctness requires the purity
observation of \cref{sec:approach}.

\textbf{GPU backends.}
\tool{} also has CUDA and Metal backends that offload DBM operations to
the GPU.  Those backends use a different parallelisation strategy (data
parallelism over the entries of all frontier DBMs simultaneously) and are
outside the scope of this paper.
